\documentclass{../../../Styles/MS-Style}
\addbibresource{../../../Assets/Citations/references.bib}

\usepackage[
    pdftitle={André Mossi CV (ES)},
    pdfauthor={André Mossi},
    pdfcreator={LaTeX y NeoVim}
]{hyperref}

\begin{document}

\begin{center}
	\authorTitle{Roberto (André) Mossi Milla}

	\contactBlock{
		Erie, Pennsylvania, 16506
	}{
		814-790-1591
	}{
		mossiroberto0392@gmail.com
	}{
		andremossi-linktree.vercel.app
	}{
		github.com/AndreM222
	}
\end{center}

\begin{Form}
	\sectionTitle{RESUMEN}
	\normalsize{
		Ingeniero de Software con una Licenciatura en Ciencias en Ingeniería de Sistemas Informáticos,
		con experiencia en sistemas de visión por computadora, investigación en inteligencia artificial,
		sistemas de simulación en tiempo real y plataformas full-stack.
		Experiencia en colaboración industrial, publicación académica y diseño integral de sistemas.
	}

	\sectionTitle{EXPERIENCIA LABORAL}

	\subSectionTitle{Gannon University}{Pensilvania, Estados Unidos}

	\subSubSectionTitle{Ingeniero de Visión por Computadora (Proyecto con Socio Industrial – NDA)}{Agosto 2024 - Mayo 2025}
	\begin{listBlock}
		\generalItem
		Arquitecté y desplegué un sistema industrial de trazabilidad en tiempo real utilizando Ultralytics YOLO
		para detectar y rastrear pallets, generando automáticamente eventos estructurados con imágenes
		y metadatos de detección.\cite{ComputerVisionVisual}

		\generalItem
		Colaboré en un equipo multidisciplinario contribuyendo en hardware, backend y componentes de software;
		asumí la responsabilidad principal del panel de monitoreo y la interfaz de visualización.

		\generalItem
		Desarrollé un dashboard web full-stack (React, JavaScript) integrado con servicios backend y TimescaleDB
		para visualizar eventos de detección, imágenes y análisis operativos en tiempo real.

		\generalItem
		Implementé herramientas administrativas dentro del dashboard para eliminar imágenes,
		corregir detecciones y curar conjuntos de datos etiquetados utilizados para el reentrenamiento iterativo
		del modelo de detección de objetos, mejorando la calidad del dataset y la confiabilidad del sistema.

		\generalItem
		Diseñé un pipeline orientado a eventos que sincroniza la captura de frames con metadatos de detección,
		habilitando flujos posteriores de análisis de anomalías y daños.

		\generalItem
		Presenté demostraciones técnicas recurrentes y revisiones del sistema a stakeholders industriales
		y evaluadores internos, asegurando financiamiento continuo e informando decisiones arquitectónicas,
		incluyendo la migración a hardware NVIDIA Jetson para mejorar rendimiento y estabilidad térmica.

		\generalItem
		Contenericé servicios distribuidos utilizando Docker e integré componentes en múltiples lenguajes
		(C++, Python, Rust) para coordinar inferencia, procesamiento de datos y comunicación del sistema.
	\end{listBlock}

	\subSubSectionTitle{Asistente de Investigación}{Agosto 2024 - Mayo 2025}
	\begin{listBlock}
		\generalItem
		Realicé investigación bajo supervisión docente en inteligencia artificial multi-agente
		simulando selección natural y comportamientos de supervivencia en entornos dinámicos.\cite{TagResearchRecommendation}

		\generalItem
		Implementé NeuroEvolution of Augmenting Topologies (NEAT) para evolucionar arquitecturas de redes neuronales
		y estrategias de comportamiento a lo largo de generaciones.

		\generalItem
		Desarrollé un entorno de simulación 3D multi-agente en Unreal Engine con generación procedural del mundo,
		interacciones basadas en física y pipelines de aprendizaje por refuerzo utilizando C++, CUDA y PyTorch.

		\generalItem
		Artículo de investigación aceptado y publicado por la American Society for Engineering Education (ASEE)
		y presentado en una conferencia Sigma Xi en Penn State University.\cite{AITagResearchPublication} \cite{AITagSigmaXIPresentation}
	\end{listBlock}

	\subSubSectionTitle{Asistente de Professor}{Agosto - Diciembre 2024}
	\begin{listBlock}
		\generalItem
		Voluntario como Asistente de Professor, apoyando a estudiantes resolviendo dudas,
		aclarando conceptos y reforzando el material del curso durante las sesiones.
	\end{listBlock}

	\subSectionTitle{Dinant}{Tegucigalpa, Honduras}

	\subSubSectionTitle{Desarrollador Full-Stack}{Junio - Julio 2023}
	\begin{listBlock}
		\generalItem
		Lideré el diseño y desarrollo integral de dos plataformas web en producción, realizando levantamiento de requisitos,
		planificación del sistema, diseño de arquitectura de bases de datos, presentaciones al cliente y despliegue final.\cite{DinantRecommendation}

		\generalItem
		Arquitecté e implementé un sistema de gestión del ciclo de vida de contratos para una empresa de venta,
		alquiler y reparación de dispositivos. Diseñé esquemas relacionales SQL para rastrear el estado de los dispositivos
		(vendido, alquilado, en reparación, devuelto), elegibilidad de servicio, filtrado por país y registros de clientes en múltiples regiones.

		\generalItem
		Desarrollé una aplicación web multi-rol (Oracle APEX, JavaScript, HTML, CSS) que permite la firma de contratos
		por invitados, captura de firma digital, generación automática de PDFs y paneles administrativos con filtrado
		avanzado por fecha, país, estado del dispositivo y estado del contrato.

		\generalItem
		Diseñé un sistema de flujo de trabajo para reparaciones con transiciones de estado, registro de daños,
		seguimiento de costos, control de tiempo y reportes diarios de actividad para garantizar responsabilidad
		y transparencia operativa.

		\generalItem
		Construí una plataforma de gestión de mantenimiento industrial que permite reportes en tiempo real
		del estado de maquinaria, flujos de mantenimiento por etapas, notificaciones automáticas, registros de reparación
		y seguimiento de uso de materiales con deducción automática de inventario y reportes para compras.

		\generalItem
		Colaboré directamente con stakeholders internacionales, realicé demostraciones iterativas,
		incorporé retroalimentación de producción y desplegué los sistemas finales utilizados por equipos operativos.
	\end{listBlock}

	\subSectionTitle{Gannon University}{Pensilvania, Estados Unidos}

	\subSubSectionTitle{Técnico Estudiantil}{Agosto - Diciembre 2022}
	\begin{listBlock}
		\generalItem
		Trabajé en el departamento de TI de la universidad, brindando asistencia en problemas tecnológicos
		a estudiantes y profesores cinco días a la semana.
	\end{listBlock}

	\sectionTitle{EDUCACIÓN}

	\subSectionTitle{Gannon University}{Pensilvania, Estados Unidos}
	\educationInfo{Licenciatura en Ingeniería de Sistemas Informáticos}{Mayo 2025\cite{GannonBSDiploma}}

	\subSectionTitle{Harvard University}{En línea}
	\educationInfo{CS50: Introducción a la Ciencia de la Computación}{Diciembre 2022\cite{HarvardCS50Certification}}

	\subSectionTitle{Extreme Networks}{En línea}
	\educationInfo{Introducción a Redes del Futuro}{Mayo 2022\cite{ExtremeCertification}}

	\sectionTitle{PUBLICACIONES}

	\begin{listBlock}
		\generalItem
		Mossi, R. A. — “Tag AI-Sandbox,” Proceedings of the ASEE Annual Conference, 2025 (revisión por pares).\cite{AITagResearchPublication}
	\end{listBlock}

	\sectionTitle{HABILIDADES ADICIONALES}

	\begin{listBlock}
		\generalItem Dominio de PowerShell, Bash y Neovim.
		\generalItem Experiencia desarrollando y desplegando aplicaciones web con React, Next.js, Vite y Vercel.
		\generalItem Nivel intermedio de C++.
		\generalItem Experiencia en automatización de compilación y scripting con CMake, Bash, PowerShell y fish.
		\generalItem Configuración y personalización de entornos Linux para profundizar en sistemas operativos
		y comportamiento de bajo nivel.
		\generalItem Fluidez escrita y oral en español e inglés.
		\generalItem Nivel intermedio escrito y oral en japonés.
	\end{listBlock}

	\sectionTitle{RESUMEN DE CALIFICACIONES}

	\begin{listBlock}
		\generalItem 3 años de experiencia en Windows PowerShell.
		\generalItem 3 años de experiencia desarrollando videojuegos en Unreal Engine.
		\generalItem 1 año de experiencia depurando software con Visual Studio, IntelliJ, Neovim y VSCode.
		\generalItem 1 año de experiencia desarrollando software en C++.
		\generalItem 5 meses de experiencia desarrollando aplicaciones web en JavaScript, HTML y CSS.
		\generalItem 3 meses de experiencia desarrollando software en Java, Python, C\# y C.
		\generalItem 2 meses de experiencia desarrollando servicios cloud a gran escala.
		\generalItem 2 meses de experiencia trabajando con Oracle y Oracle Apex.
		\generalItem 2 meses de experiencia desarrollando sitios web en NodeJS, React y Next.js.
	\end{listBlock}

	\sectionTitle{PREMIOS Y RECONOCIMIENTOS}

	\begin{listBlock}
		\awardItem{Premio al Desarrollo Escolar}{Reconocimiento por el desarrollo destacado de un videojuego utilizando Blueprints en Unreal Engine 4}{2020\cite{SoftwareEngineer2020Award}}

		\awardItem{Premio al Desarrollo Escolar}{Reconocimiento por el desarrollo destacado de un videojuego utilizando C++ y Unreal Engine 4}{2016\cite{SoftwareEngineer2016Award}}
	\end{listBlock}

	\bibliographyBlock[REFERENCIAS]

\end{Form}
\end{document}
